\section*{Cinq principes holographiques de construction et de conservation}
\addcontentsline{toc}{section}{Cinq principes holographiques de construction et de conservation}

Les \emph{Éléments} offrent un antécédent précis pour présenter une
architecture mathématique au moyen d’opérations déclarées et de conséquences
démontrées. Le livre I distingue les définitions, les postulats et les notions communes.
Ses trois premiers postulats permettent de tracer un segment entre des points,
de prolonger un segment en ligne droite et de décrire un cercle de centre et de rayon
donnés ; le quatrième établit l’égalité des angles droits. Le cinquième
formule une condition de rencontre : deux droites coupées par une transversale
se rencontrent, lorsqu’elles sont prolongées, du côté où les angles intérieurs ont une somme
inférieure à deux angles droits. Ce sont des énoncés sur les objets et les constructions de
cette géométrie.\footnote{Euclide, \emph{Éléments}, livre I, postulats 1--5.
Texte grec de Heiberg et traduction de Richard Fitzpatrick, p.~7 :
\url{https://farside.ph.utexas.edu/books/Euclid/Elements.pdf}. On peut aussi
les consulter séparément dans l’édition de David E. Joyce :
\url{https://mathcs.clarku.edu/~djoyce/elements/bookI/bookI.html}.}

L'extrême et moyenne raison apparaît dans le livre VI comme une relation entre un tout et ses parties : le tout est au plus grand segment comme celui-ci est au plus petit. Dans une traduction algébrique moderne, pour des longueurs positives \(u>v\), on écrit \((u+v)/u=u/v\) ; par conséquent, le rapport \(r=u/v\) satisfait \(r^2=r+1\). Cette écriture explique la proportion transmise sans faire du langage algébrique moderne le point de départ d'Euclide. Dans HMT, la coordonnée d'auto-échelle s'obtient à partir d'une incidence produite par le calendrier opératoire ; sa reconnaissance au moyen de la même équation appartient à une étape ultérieure.\footnote{Euclide, \emph{Éléments}, livre VI,
définition 3, dans l’édition de Joyce :
\url{https://mathcs.clarku.edu/~djoyce/elements/bookVI/defVI3.html}.
L’édition de Fitzpatrick imprime ce même énoncé comme définition 2
du livre VI, p.~156. Les deux numérotations éditoriales sont conservées ici.}

La comparaison qui suit se situe dans l’organisation de la construction.
HMT déclare un support arithmétique, deux feuillets reliés, une orientation
tritique et des opérations qui prolongent les états en conservant leurs registres. Un
cylindre représente l’information fixée par un préfixe ; une droite
euclidienne est un autre objet, avec une autre définition et d’autres opérations. Les
cinq principes de cette section expriment des relations effectives entre les
objets holographiques. Leur nombre offre une exposition reconnaissable ; il
n’établit ni correspondance terme à terme avec les cinq postulats,
ni réfutation du cinquième, ni remplacement des axiomes géométriques.

Ces principes ont un statut \emph{organisateur} : ils réunissent des conditions et des
théorèmes que le corps de l’article construit et prouve. Chacun présente d’abord
son sens, puis ses domaines et une identité qui permet de le vérifier.
La suite commune est
\[
 \mathrm{APP}\longrightarrow\mathrm{TRIT}\longrightarrow\mathrm{TPK}
 \longrightarrow\text{état enrichi}
 \longrightarrow\text{structure discrète conjointe du continu}.
\]
Les réalisations numériques, linéaires et géométriques lisent ensuite cette structure. Cette position des réalisations maintient visibles l'origine des coefficients et l'information que conserve chaque lecture.

\subsection*{Principe I. Unité relationnelle de la génération}
\hypertarget{prinh:unidad}{}

\emph{Chaque lecture est construite par des opérations typées sur un
même support, en conservant les relations qui permettent de la continuer.} Une
case APP admet simultanément une lecture additive et une lecture multiplicative.
Le passage à un chiffre visible enregistre une partie des deux opérations ; le
quotient, l’orientation et le feuillet conservent les conditions de son
prolongement. L’unité consiste dans cette provenance relationnelle, qui permet de
suivre une même exécution à travers ses différentes réalisations.

Le domaine initial est \(D_9^2\), avec \(D_9=\{1,\ldots,9\}\). Sur celui-ci
agissent \(S(i,j)=i+j\) et \(P(i,j)=ij\). Pour un entier \(n\), le système de coordonnées
positif du résidu et du quotient est donné par
\[
 \rho_9(n)=1+((n-1)\bmod9),\qquad
 q_9(n)=\frac{n-\rho_9(n)}9,\qquad
 n=\rho_9(n)+9q_9(n).
\]
L’évaluation relevée d’APP conserve
\(\bigl((\rho_9(S),q_9(S)),(\rho_9(P),q_9(P))\bigr)\).
L’identité reconstruit les deux lectures entières ; par exemple, elle distingue
les valeurs 10 et 19, qui ont le même résidu positif et des quotients
différents. Le résultat et sa preuve sont développés dans
\eqref{eq:residuo-cociente} et \eqref{eq:app-levantada}. La représentation
par résidus \(0,\ldots,8\), utilisée dans d’autres systèmes de coordonnées de l’article, conserve
son quotient correspondant : changer de représentants exige de changer les deux
composantes de la division.

TRIT ajoute la représentation orientée de \(\mathbb F_3\) par \(\{-1,0,+1\}\), avec sa division exacte \(n=r+3c\). Le sélecteur opératoire \(M_{\rm ph}\) prend les valeurs \(+1,-1,0\) dans les classes de phase \(\{1,4,7\}\), \(\{2,5,8\}\), \(\{3,6,9\}\), respectivement. Le transport cardinal a un autre domaine et une autre fonction : il déplace les curseurs sur \(X_9=(\mathbb Z/9\mathbb Z)^2\). Dans l'état enrichi \(\widehat X\), TPK compose sélection, transport et mise à jour de mémoire,
\[
 U_t=\operatorname{Mem}_t\circ\operatorname{Tra}_t
                    \circ\operatorname{Sel}_t.
\]
La définition \eqref{eq:actualizacion} et le développement de
\cref{tpk-int:estado,tpk-int:seleccion} spécifient la phase, les curseurs, les feuillets,
les directions, les résidus, les quotients et les registres. Le signe de phase et la direction
cardinale se composent en conservant leurs types distincts.

L’auto-échelle offre un exemple de génération que l’on peut suivre sans partir
de sa valeur finale. Le calendrier \((+1,-1,0)\) produit, avec la règle de
changement et de conservation de feuillet, le cycle modal
\((\Sigma,\Pi,\Pi)\). Ses trois arêtes déterminent
\[
 F_{\rm av}=\begin{pmatrix}0&1\\1&1\end{pmatrix}.
\]
Dans sa réalisation linéaire ultérieure, le rayon propre positif normalisé comme \((1,x)^{\mathsf T}\) satisfait \(x=\lambda\) et \(1+x=\lambda x\). On obtient ainsi \(x^2-x-1=0\) ; la racine positive est reconnue comme \(\varphi\). Le théorème qui suit \eqref{eq:fav} prouve cette identification après la construction des entrées de la matrice. La relation avec la proportion euclidienne est exacte dans cette lecture d'auto-échelle et conserve la différence entre leurs deux procédures d'origine.

\subsection*{Principe II. Prolongement orienté et mémoire de composition}
\hypertarget{prinh:prolongacion}{}

\emph{Un prolongement admissible ajoute un événement à l’état et
transporte la mémoire antérieure au moyen de la même opération.} L’histoire
permet de distinguer des parcours qui se terminent à une position ou à une phase
identiques. Ce principe fournit une manière précise de parler de
continuation : on conserve le préfixe effectivement parcouru et on ajoute
une émission permise par le sélecteur et le régime présents.

Soit \(X_k\) l’ensemble des histoires admissibles de longueur \(k\),
construites à partir des germes du noyau. Si \(\varepsilon\) est admissible
après \(h\), le prolongement \(h\varepsilon\) appartient à \(X_{k+1}\)
et la troncature satisfait \(\rho_k(h\varepsilon)=h\). On retient
l’étiquette d’émission et les feuillets antérieurs, même lorsque deux émissions
ont la même lecture finale. Dans un système de coordonnées de mémoire sur un module
\(M\), la mise à jour a la forme
\(U_\varepsilon(m)=C_\varepsilon m+\kappa_\varepsilon\), avec
\(C_\varepsilon\in\operatorname{End}(M)\) et
\(\kappa_\varepsilon\in M\). Si deux tronçons sont composables dans ce
système de coordonnées, alors
\[
 C_{\beta\alpha}=C_\beta C_\alpha,\qquad
 \kappa_{\beta\alpha}=C_\beta\kappa_\alpha+\kappa_\beta,
 \qquad U_{\beta\alpha}=U_\beta U_\alpha.
\]
La preuve de \cref{iv:cont:concatenacion} consiste à substituer la première mise à jour dans la seconde. La mémoire est additionnée après le transport qui lui correspond ; l'associativité rend le résultat indépendant de la subdivision du trajet.

La différence entre phase et mémoire apparaît déjà dans une lecture élémentaire.
Pour \((w,r)\in\mathbb Z\times C_9\), avec
\(\bar r\in\{0,\ldots,8\}\), on a
\[
 \sigma_9(w,r)=
 \left(w+\left\lfloor\frac{\bar r+1}{9}\right\rfloor,
       r+1\pmod9\right),\qquad
 \sigma_9^9(w,r)=(w+1,r).
\]
La seconde égalité compte exactement un passage de 8 à 0. Sa première coordonnée enregistre l'avancée que la seconde ne distingue plus. Dans le calendrier de douze fenêtres, la suite exacte \(0\to C_{12}\to C_{108}\to C_9\to0\) conserve la retenue entre fenêtres. Elle ne se scinde pas : \(C_{108}\) est d'exposant 108 et \(C_{12}\times C_9\) d'exposant 36. La preuve et la loi de la retenue figurent dans \cref{prop:k-extension}.

Le retour de phase acquiert ainsi une signification opératoire : il permet de
comparer des lectures effectuées à une même position du calendrier tandis que
l’état conserve ce qui s’est passé entre elles. L’holonomie \(\Gamma_9\)
s’obtient à partir de la dynamique enrichie ; \(K_{\rm ph}\) est sa représentation
réduite ultérieure. Le registre dodécaphasique entier \(K\), qui utilise
en outre l’orientation et les incidences, a également sa construction et son domaine
propres. L’égalité d’une phase, d’un résidu ou d’une coordonnée exprimée
n’identifie pas à elle seule les histoires. C’est l’information de départ
qui est conservée lors du passage au raffinement.

\subsection*{Principe III. Compatibilité du raffinement et réalisation}
\hypertarget{prinh:refinamiento}{}

\emph{Raffiner consiste à ajouter une résolution compatible avec le préfixe,
en maintenant les relations de troncature et les masses de ses prolongements.}
La lecture à une échelle plus fine précise ce qui s’est produit à l’intérieur d’un
cylindre ; lorsqu’on oublie ce détail, on récupère la description précédente.
La réalisation d’une valeur surgit lorsque les cylindres de sa lecture se
contractent de manière compatible. L’histoire qui détermine ces cylindres
demeure dans le domaine de la construction.

Dans l’arbre de \cref{iv:cont:medida}, l’ensemble initial est fini et non
vide, et chaque histoire a un nombre fini et positif de successeurs
admissibles. Une masse initiale positive de somme un et des probabilités locales
\(p(\varepsilon\mid h)>0\), de somme un en chaque prédécesseur, produisent
\[
 \mu(h\varepsilon)=\mu(h)p(\varepsilon\mid h),\qquad
 \sum_{\rho_k(y)=h}\mu(y)=\mu(h).
\]
Par conséquent, dans la réalisation ultérieure
\(H_k=L^2(X_k,\mu_k)\), le relèvement
\((J_kf)(y)=f(\rho_k(y))\) est isométrique. Son adjoint est l’espérance
conditionnelle
\[
 (E_kg)(h)=\sum_{\rho_k(y)=h}
                  \frac{\mu(y)}{\mu(h)}g(y),
 \qquad E_kJ_k=I.
\]
Regrouper la norme au carré par prédécesseurs démontre l’isométrie et regrouper le
produit scalaire démontre la formule de l’adjoint
(\cref{iv:cont:esperanza}). Le relèvement d’une lecture et son retour
au niveau précédent forment ainsi un couple d’opérations explicites.

Les hypothèses de ramification finie et de prolongement non vide fournissent
des histoires infinies compatibles. Leurs cylindres reçoivent les masses précédentes,
qui déterminent une mesure sur la tribu cylindrique. Le choix de
\(p\) appartient au lecteur déclaré : le cas de successeurs équiprobables donne
\(p=1/d(h)\), sans imposer cette mesure à toute autre lecture. Cette
construction conserve les multiplicités des histoires, y compris celles qui
partagent une valeur réalisée. La prolongeabilité se réfère à l’arbre
admissible construit ; une restriction supplémentaire des routes doit conserver
expressément cette propriété.

Pour la lecture numérique régionale, chaque bloc
\(b\in\mathbb F_3^6\) possède un indice entier
\(\nu(b)=\sum_{j=1}^6 b_j3^{6-j}\), compris entre 0 et 728.
Le lecteur fixe \(N_0=m\) et produit

\[
 N_{n+1}=729N_n+\nu(b_{n+1}),\qquad
 I_n=\left[\frac{N_n}{729^n},\frac{N_n+1}{729^n}\right].
\]
On a \(I_{n+1}\subseteq I_n\) et
\(\operatorname{diam}(I_n)=729^{-n}\). Les extrémités gauches forment
une suite rationnelle de Cauchy ; sa classe définit d’abord la réalisation
interne de la lecture. Identifiée ultérieurement à la complétion
usuelle des rationnels, elle correspond à l’unique point des intervalles
fermés emboîtés. Le théorème postérieur à \eqref{eq:cilindro} démontre
cette assertion et traite séparément les conventions de frontière. Pour
les caractères régionaux, le lecteur effectif de chaque préfixe décimal
utilise les encadrements calculables et la séparation démontrée par rapport à
leurs frontières, comme le précise \cref{gen:prefijos-regionales}. Cette profondeur
arbitraire conserve le quantificateur « pour chaque longueur finie » ;
l’existence d’une limite par compatibilité et contraction est une assertion
distincte et ne fournit pas, à elle seule, un algorithme pour toute histoire
abstraite. Les histoires qui représentent le même point conservent leurs
registres distincts.

\subsection*{Principe IV. Double lecture corrélative et transport de repère
}
\hypertarget{prinh:doble}{}

\emph{Les lectures numérique et incidentielle procèdent du même état
enrichi et conservent leurs compatibilités lors du prolongement et du changement de
repère.} La première précise une réalisation au moyen d’intervalles de plus en plus
étroits. La seconde retient les relations entre les positions, les signes et les
frontières. Le lien holographique consiste en leurs applications communes d’origine et de
transport, qui permettent de comparer quelle information chacune exprime.

Soit \(X_n^{\rm cal}\) le domaine des préfixes admissibles qui incluent leur repère initial et les transports successifs de repère. La lecture arithmétique utilise les blocs de ce préfixe ; la lecture d'incidence utilise les mots \(w_j\in\mathbb F_3^6\) et les repères \(f_j\). Pour une matrice \(A\in M_6(\mathbb F_3)\), le lecteur linéaire
\[
 \gamma_A:\mathbb F_3^6\longrightarrow\mathbb F_3^{12},
 \qquad \gamma_A(w)=(w,wA)
\]
agit sur des mots écrits sous forme de lignes. Sa spécialisation dans la matrice
Paley--Witt construite dans \cref{exc:paley-codigo} produit le code
ternaire et ses incidences. La matrice est obtenue dans cette construction
finie ; un nombre réel terminal ne sélectionne pas ses entrées. L’ordre
\(w_6\to w_{12}\to w_{18}\to w_{24}\to w_{30}\to R_{36}\to G_9\)
maintient les transformations et les mémoires utilisées avant chaque lecture.


La compatibilité peut se vérifier sur l'opérateur lui-même. Pour \(f\in\operatorname{GL}_6(\mathbb F_3)\), soit \(A^f=fAf^{-1}\). Alors
\[
 \gamma_{A^f}(wf^{-1})
      =\gamma_A(w)(f^{-1}\oplus f^{-1}),
\]
puisque les deux côtés sont \((wf^{-1},wAf^{-1})\). En particulier, la lecture du préfixe
\[
 Q_n(x)=\bigl(f_j,
       \gamma_{f_jA_Wf_j^{-1}}(w_j)\bigr)_{0\leq j<n}
\]
satisfait \(r_{n,m}Q_n=Q_m p_{n,m}\), où \(p_{n,m}\) tronque les histoires et \(r_{n,m}\) tronque les registres. La causalité de la mise à jour des repères garantit que les deux procédures retiennent exactement les mêmes termes. Le résultat \cref{exc:limite-inverso} construit son passage à la limite. La covariance linéaire est valable pour le changement de repère écrit ; conserver en outre un support coordonné ou une métrique exige de transporter cette structure, ou de restreindre les changements de repère au groupe qui la préserve.

La construction corrélative du continu étend cette compatibilité au
transport par feuillets, au bilan orienté, à l’incidence ternaire, au
complexe rectangulaire de ports et à la mémoire cellulaire avec ses remplissages et
ses droites déterminantes. Ces aspects se construisent sur les mêmes
histoires, non sur cinq domaines choisis après coup. Les identités de
concaténation, de naturalité et de bord assurent leur assemblage dans
\cref{iv:rectangular-natural,iv:cont:memoria-natural,iv:cont:ensamblaje,iv:cont:mapa-correlativo}.
Les cinq principes d’exposition de cette ouverture ne constituent pas une nouvelle
partition de ces cinq constructions consubstantielles. La récupération
d’un mot au moyen de la première projection de \(\gamma_A\), quant à elle,
concerne ce mot : la récupération de l’histoire enrichie
complète exige de conserver les autres registres qui la distinguent.


\subsection*{Principe V. Conservation et récupération du registre complet}
\hypertarget{prinh:conservacion}{}

\emph{L’information récupérable est conservée conjointement dans l’état
qui se poursuit, dans les registres produits et dans leurs corrélations.} Une
lecture visible peut se contracter tandis que le registre complet conserve une
représentation isométrique de l’état de départ. Ce principe permet de
formuler une réversibilité vérifiable : on déclare ce qui est stocké, dans quel
espace cela réside et par quel opérateur cela est reconstruit.

Le transport d'histoires relevées de \cref{iv:cont:transporte} agit entre les espaces \(H_k=\ell^2(\widehat X_k)\) au moyen de
\[
 \mathcal U_ke_{\widehat h}=
 \sum_{\varepsilon\ {
m admisible}}
 \sqrt{p(\varepsilon\mid\widehat h)}\,
 e_{\widehat h\varepsilon},\qquad
 \mathcal U_k^*\mathcal U_k=I.
\]
La seconde égalité procède de la normalisation locale et du fait que des prédécesseurs
distincts ont des successeurs aux étiquettes distinctes. Si \(P_k\) lit le feuillet
\(\Sigma\) et \(Q_k=I-P_k\) le feuillet \(\Pi\), les opérations
\[
 \begin{aligned}
 A_k&=P_{k+1}\mathcal U_kP_k+Q_{k+1}\mathcal U_kQ_k,\\
 \eta_k&=Q_{k+1}\mathcal U_kP_k+P_{k+1}\mathcal U_kQ_k
 \end{aligned}
 \qquad\text{satisfont}\qquad A_k^*A_k+\eta_k^*\eta_k=I.
\]
L'identité de Gram de \cref{iv:cont:gram} se démontre en annulant les produits de projecteurs complémentaires. Elle distingue le fait de rester dans un feuillet de celui de changer de feuillet ; la valeur tritique zéro reste une orientation et ne devient pas un troisième feuillet.

Écrivons \(F_0=I\), \(F_{n+1}=A_nF_n\). Le télescopage de
l’identité précédente donne, pour tout horizon \(N\),
\[
 \|x\|^2=\|F_Nx\|^2+
                 \sum_{n=0}^{N-1}\|\eta_nF_nx\|^2.
\]
Le premier terme est le résidu terminal qui se poursuit encore ; les autres
sont les sorties enregistrées. Le théorème \cref{iv:cont:terminal} démontre
que \(F_N^*F_N\) converge fortement vers un opérateur positif
\(T_\infty\). On obtient donc le registre complet

\[
 \mathcal V:H_0\longrightarrow H_0\oplus
                     \bigoplus_{n\geq0}H_{n+1},\qquad
 \mathcal Vx=\bigl(T_\infty^{1/2}x,
                         (\eta_nF_nx)_{n\geq0}\bigr),
 \qquad\mathcal V^*\mathcal V=I.
\]
En effet, la somme des carrés des normes est \(\|x\|^2\) ; la polarisation donne la dernière identité. Son image \(\mathcal M\) est fermée et l'inverse sur celle-ci est \(\mathcal V^*\). Le terminal est conservé jusqu'à ce que son extinction soit démontrée ; \(T_\infty\) n'a pas besoin d'être un projecteur et la convergence de \(F_Nx\) n'est pas supposée.

La conservation atteint également les opérateurs qui lisent cette information.
Pour un opérateur borné \(B\) sur \(H_0\), le transport
\(B^{\mathcal V}=\mathcal VB\mathcal V^*|_{\mathcal M}\) conserve
les produits, les adjoints et les produits scalaires :
\[
 (B_1B_2)^{\mathcal V}=B_1^{\mathcal V}B_2^{\mathcal V},\qquad
 \langle\mathcal Vx,B^{\mathcal V}\mathcal Vy\rangle
                       =\langle x,By\rangle.
\]
Insérer \(\mathcal V^*\mathcal V=I\) prouve les deux formules. Dans le
lecteur nonadique de \cref{lib01:isometria}, les mêmes identités se
réalisent avec \(T_n=(8I+C_n)/9\) et
\(D_n=\sqrt8(C_n-I)/9\), pour des unitaires \(C_n\) du porteur
déclaré. La composition maintient son statut de lecteur et n’identifie pas
automatiquement chaque \(C_n\) à l’holonomie génératrice \(\Gamma_9\).
La récupération préserve ainsi les corrélations du registre et les
observables transportés, dans l’image précise où elle agit.

Le registre \(K\) permet de voir cette distinction dans un objet fini. Son
codage \(\kappa(K)\), en base \(B=1000\) et de longueur douze, est

\[
 \kappa(K)=\frac{N(K)}{B^{12}-1},\qquad
 N(K)=\sum_{j=1}^{12}K_jB^{12-j},\qquad
 K\in\{0,\ldots,999\}^{12}.
\]
Multiplier par \(B^{12}-1\) et effectuer douze divisions euclidiennes
retrouve exactement le registre, y compris ses blocs initiaux nuls.
La proposition \cref{lib01:inversion-kappa} prouve en outre qu’une erreur
absolue inférieure à \(1/[2(B^{12}-1)]\) permet de retrouver le même entier
par arrondi. La réversibilité porte sur le registre dans son système de coordonnées et
avec ses métadonnées ; elle ne l’identifie pas à toute l’histoire TPK. La
\emph{constante holographique de couplage arithmético-géométrique} désigne
dans cet article l’objet \(K\) typé avec ses lecteurs et ses relations. Sa
représentation procède de l’état enrichi et participe ensuite à la
clôture d’alpha : elle n’est pas utilisée comme entrée génératrice de \(\pi\), et
sa dénomination ne transforme pas chaque lecture scalaire en invariant universel.

\subsection*{Une lecture conjointe des cinq principes
}

La séquence offre un critère positif pour parcourir l’article. L’unité
relationnelle fixe ce sur quoi on opère ; le prolongement conserve ce qui a été parcouru ; le
raffinement organise la résolution ; la double lecture maintient ensemble valeur
et incidence ; la récupération identifie les conditions exactes permettant de
reconstruire le registre et ses observables. Les démonstrations s’enchaînent
parce qu’elles agissent sur des objets de provenance commune et respectent leurs applications de
transport. Là où une lecture oublie des informations, son domaine et sa
fibre de perte restent définis.


L’extrême et moyenne raison historique montre comment une relation entre le tout
et ses parties peut déterminer une proportion. La construction holographique
développée ici étudie, en outre, des relations conservées lorsque changent la
résolution, le repère et la distribution de l’information entre continuité et
mémoire. L’équation d’auto-échelle établit un lien mathématique précis
avec cette proportion ; les identités de concaténation, de raffinement et
d’isométrie fournissent le contenu supplémentaire que cet article développe.
La comparaison historique conserve ainsi sa fonction d’orientation, tandis que
chaque conclusion holographique repose sur ses opérations et ses preuves propres.

\subsection*{Sources primaires de la comparaison historique
}

\begin{enumerate}[label=\arabic*.]

 \item Euclide, \emph{Éléments}, livre I, postulats 1--5. Texte
 grec établi par J. L. Heiberg, avec traduction de Richard Fitzpatrick,
 édition corrigée de 2008, p.~7 ; livre VI, définition 2 dans cette édition,
 p.~156. \url{https://farside.ph.utexas.edu/books/Euclid/Elements.pdf}.

 \item Euclide, \emph{Éléments}, édition numérique de David E. Joyce,
 Clark University : livre I, postulats 1--5,
 \url{https://mathcs.clarku.edu/~djoyce/elements/bookI/bookI.html} ;
 livre VI, définition 3,
 \url{https://mathcs.clarku.edu/~djoyce/elements/bookVI/defVI3.html}.
 La différence de numérotation du livre VI est conservée pour faciliter la comparaison entre éditions.
\end{enumerate}
